\input{../tex-inputs/latex-leseansicht-vorspann}

\section[Gustav Schwarzkopf an Arthur Schnitzler, 18. 9. 1900]{L04316 Gustav Schwarzkopf an Arthur Schnitzler, 18. 9. 1900}
\nopagebreak\mylabel{L04316v}
\rehead{ }\normalsize\beginnumbering\briefempfaengerindex{Schnitzler, Arthur@\textsc{Schnitzler, Arthur}!zzzHiller, Max@\emph{von Max Hiller}!1900-09-181@{18. 9. 1900}|(be}\briefempfaengerindex{Schnitzler, Arthur@\textsc{Schnitzler, Arthur}!zzzHiller, Ida@\emph{von Ida Hiller}!1900-09-181@{18. 9. 1900}|(be}\briefempfaengerindex{Schnitzler, Arthur@\textsc{Schnitzler, Arthur}!zzzFoges, Max@\emph{von Max Foges}!1900-09-181@{18. 9. 1900}|(be}\briefempfaengerindex{Schnitzler, Arthur@\textsc{Schnitzler, Arthur}!zzzSchwarzkopf, Gustav@\emph{von Gustav Schwarzkopf}!1900-09-181@{18. 9. 1900}|(be}
\toendnotes[C]{\smallbreak\pagebreak[2]}
\correspDesc{Versand  durch Gustav Schwarzkopf, Max Foges, Ida Hiller, Max Hiller am 18. 9. 1900 in Puchberg am Schneeberg
\newline{}Erhalt  durch Arthur Schnitzler im Zeitraum [19. 9. 1900 – 23. 9. 1900?] in Wien}\toendnotes[C]{\smallbreak}
\Standort{CUL, Schnitzler, B 96a.}
\physDesc{Bildpostkarte, 113 Zeichen
\newline{}Handschrift Gustav Schwarzkopf: Bleistift, deutsche Kurrent
\newline{}Handschrift Max Hiller: Bleistift
\newline{}Handschrift Ida Hiller: Bleistift
\newline{}Handschrift Max Foges: Bleistift}\pstart{}{\pb}Herrn \textsc{D\textsuperscript{r} Arthur Schnitzler}\pend{}\pstart{}\textsc{Wien\oindex{Wien@\textbf{Wien}, \emph{Verwaltungsgebiet}|pw}}.\pend{}\pstart{}IX. Frankgaſſe 1\oindex{Wien@\textbf{Wien}!IX., Alsergrund@\textbf{IX., Alsergrund}!Frankgasse 1@\textbf{Frankgasse 1}, \emph{Wohngebäude}|pw}\pend{}{\bigskip}
\pstart
           \noindent{}\centering{}{\pb}\textcolor{gray}{\textbf{Hôtel Hochschneeberg}}\oindex{Berghaus Hochschneeberg@\textbf{Berghaus Hochschneeberg}, \emph{Hotel}|pw}\pend
           \vspace{1em}
\pstart
           \noindent{}{\pb}Herzlichen Gruß\pend
           
\pstart
           \spacefill\mbox{Gustav Schwarzkopf}{\\[\baselineskip]}\spacefill\mbox{{[}hs. Foges:{]} Foges}{\\[\baselineskip]}\spacefill\mbox{{[}hs. Hiller:{]} MaxHiller}{\\[\baselineskip]}\spacefill\mbox{{[}hs. Hiller:{]} Ida Hiller}\pend
           \leftskip=0em{}
\pstart
           18. 9. 1900\pend
           \selectlanguage{ngerman}\endnumbering\briefempfaengerindex{Schnitzler, Arthur@\textsc{Schnitzler, Arthur}!zzzHiller, Max@\emph{von Max Hiller}!1900-09-181@{18. 9. 1900}|)be}\briefempfaengerindex{Schnitzler, Arthur@\textsc{Schnitzler, Arthur}!zzzHiller, Ida@\emph{von Ida Hiller}!1900-09-181@{18. 9. 1900}|)be}\briefempfaengerindex{Schnitzler, Arthur@\textsc{Schnitzler, Arthur}!zzzFoges, Max@\emph{von Max Foges}!1900-09-181@{18. 9. 1900}|)be}\briefempfaengerindex{Schnitzler, Arthur@\textsc{Schnitzler, Arthur}!zzzSchwarzkopf, Gustav@\emph{von Gustav Schwarzkopf}!1900-09-181@{18. 9. 1900}|)be}\mylabel{L04316h}
\begin{anhang}
\end{anhang}\newcommand{\dateiname}{L04316}\newcommand{\titel}{Gustav Schwarzkopf an Arthur Schnitzler, 18. 9. 1900}\newcommand{\editorInnen}{Herausgegeben von Jahnke, SelmaMüller, Martin Anton}\input{../tex-inputs/latex-leseansicht-abspann}
